% Part: sets-functions-relations
% Chapter: relations
% Section: reflections
%
\documentclass[../../../include/open-logic-section]{subfiles}

\begin{document}

\olfileid{sfr}{rel}{ref}
\olsection{तात्त्विक चिंतन}

\olref[set]{sec} मध्ये आपण संबंधांची व्याख्या काही विशिष्ट संच म्हणून
केली. येथे थांबून एक छोटासा तात्त्विक प्रश्न विचारू: अशी व्याख्या नेमके
काय \emph{करते}? काही सत्तामीमांसक एकरूपतेची तथ्ये आपण \emph{शोधून काढली},
असा दावा आपल्याला करायचा आहे का, याबद्दल खूपच शंका आहे. उदाहरणार्थ,
$\Nat$ वरील क्रमसंबंध हा \olref[set]{sec} मध्ये व्याख्यिलेला
$R=  \Setabs{\tuple{n,m}}{n, m \in \Nat\text{ आणि }
n < m}$ संचच \emph{निघाला}, असे म्हणायचे का? अशी शंका घेण्याची तीन
कारणे पुढीलप्रमाणे आहेत.

पहिले: \olref[set][pai]{wienerkuratowski} मध्ये आपण
$\tuple{a, b} = \{\{a\}, \{a, b\}\}$ अशी व्याख्या केली. त्याऐवजी
$\lVert a,b\rVert = \{\{b\}, \{a, b\}\} = \tuple{b,a}$ ही व्याख्या
विचारात घ्या. $a \neq b$ असेल, तर $\tuple{a, b} \neq \lVert a,b\rVert$.
पण क्रमित जोडीची व्याख्या $\tuple{a,b}$ ऐवजी $\lVert a,b\rVert$ अशीही
तितक्याच योग्य रीतीने करता आली असती. दोन्ही व्याख्या तितक्याच उपयोगी
ठरल्या असत्या. त्यामुळे नैसर्गिक संख्यांवरील क्रमसंबंध ``असण्यासाठी''
आता दोन तितकेच योग्य उमेदवार आहेत:
\begin{align*}
		R &= \Setabs{\tuple{n,m}}{n, m \in \Nat \text{ आणि }n < m}\\
		S &= \Setabs{\lVert n,m\rVert}{n, m \in \Nat \text{ आणि }n < m}.
\end{align*}
विस्तारात्मकतेनुसार $R \neq S$ असल्याने ते \emph{दोन्ही} $\Nat$ वरील
क्रमसंबंधाशी एकरूप असू शकत नाहीत. पण प्रत्यक्षात $S$ ऐवजी $R$ (किंवा
उलट) \emph{हाच} क्रमसंबंध आहे, असा दावा करणे मनमानीचे आणि म्हणून काहीसे
अडचणीचे ठरेल. (संचसैद्धान्तिक न्यूनीकरणवादाविरुद्ध बेनासेराफ यांनी
\citeyear{Benacerraf1965} मध्ये प्रसिद्ध केलेल्या युक्तिवादाचे हे अतिशय
सोपे उदाहरण आहे. आपण त्याकडे आणखी काही वेळा परत येऊ.)

दुसरे: \emph{प्रत्येक} संबंध एखाद्या संचाशी एकरूप मानला पाहिजे असे आपल्याला
वाटत असेल, तर संचसदस्यत्वाचा $\in$ हा संबंधसुद्धा एक विशिष्ट संच असला
पाहिजे. तो $\Setabs{\tuple{x,y}}{x \in y}$ हा संच असावा लागेल. पण हा
संच अस्तित्वात आहे का? रसेलची विरोधापत्ती लक्षात घेता, असा संच अस्तित्वात
आहे हा दावा सहज स्वीकारता येत नाही. खरे तर,
\oliflabeldef{cumul:::part}{\olref[cumul][][]{part} मध्ये आपण विकसित
करणारा संच सिद्धांत या संचाचे अस्तित्व \emph{नाकारील}.\footnote{पुढील
चर्चेची थोडी झलक अशी. व्याघात काढण्यासाठी $I =
\Setabs{\tuple{x,y}}{x \in y}$ अस्तित्वात आहे असे समजा. मग
$\bigcup \bigcup I$ हा सार्वत्रिक संच होतो; हे
\olref[sfr][z][sep]{thm:NoUniversalSet} शी विसंगत आहे.}}{संच सिद्धांत
काटेकोरपणे स्वयंसिद्धकांवर आधारित सिद्धांत म्हणून विकसित करता येतो,
आणि तो सिद्धांत या संचाचे अस्तित्व नाकारतो.}
म्हणून काही संबंधांना संच म्हणून वागवता येत असले, तरी संचसदस्यत्वाचा
संबंध अपवादात्मक प्रकरण म्हणून घ्यावा लागेल.

तिसरे: संबंध संचांशी ``एकरूप'' मानताना, $\tuple{x,y} \in R$ यासाठी
$Rxy$ असे लिहिण्याची मुभा आपण घेतली. सदस्यत्वाचा ``$\in$'' हा संबंध
विधेय \emph{म्हणून} वापरला, तर हे योग्य आहे. पण ``$\in$'' हे एका
विशिष्ट प्रकारच्या संचाचे नाव आहे असे मानले, तर ``$\tuple{x,y} \in R$''
या व्यंजकात संचांचा निर्देश करणारी फक्त तीन एकवाची पदे राहतात:
``$\tuple{x,y}$'', ``$\in$'' आणि ``$R$''. अशा नावांच्या यादीतून विधान
व्यक्त होऊ शकत नाही; ``तो कप लेखणीधारक ते टेबल'' या निरर्थक शब्दमालेइतकीच
ती निष्फळ आहे. म्हणून पुन्हा, काही संबंधांना संच म्हणून वागवता येत असले,
तरी संचसदस्यत्वाचा संबंध अपवादात्मक प्रकरण असला पाहिजे. (यात फ्रेगे यांच्या
\emph{घोडा} संकल्पनेच्या विरोधापत्तीचे एक सोपे रूप आणि विट्गेन्स्टाइन यांनी
रसेल यांच्याविरुद्ध मांडलेला एक प्रसिद्ध आक्षेप एकत्र आले आहेत.)

मग यातून काय निष्पन्न होते? संख्यांवरील संबंध हे संच आहेत असे म्हणण्यात
काही \emph{चूक} नाही. ते म्हणण्यामागचा हेतू मात्र नीट समजून घ्यावा लागेल.
आपण सत्तामीमांसक एकरूपतेचे तथ्य सांगत नाही. फक्त एवढेच नोंदवत आहोत की,
काही संदर्भांत काही संबंधांना काही विशिष्ट संच \emph{म्हणून वागवता}
येते (आणि आपण तसे वागवणार आहोत).

\end{document}
